\chapter{Два огляди сенсорного злиття: embodied AI і event odometry}
\label{chap:dual-sensor-fusion}

\sourcebox{Source-preserved модулі: \sourcepath{sensor_fusion_2506.19769/} і \sourcepath{sensor_fusion_2410.15480/}. Другий модуль у цій сесії скомпільовано як окремий IEEE-style PDF.}

Після останнього auxiliary local run payload сенсорне злиття стало окремою сильною віссю корпусу. Є два різні, але взаємодоповнювальні матеріали.

\section{Fusion taxonomy}

Для прикладної математики корисно розділяти рівні злиття:
\begin{itemize}[leftmargin=*]
\item \textbf{data-level fusion}: об'єднання сирих або майже сирих вимірювань;
\item \textbf{feature-level fusion}: об'єднання ознак, embeddings, keypoints, point/voxel/BEV representations;
\item \textbf{decision-level fusion}: об'єднання уже готових рішень або гіпотез;
\item \textbf{temporal fusion}: послідовне накопичення інформації в часі;
\item \textbf{multi-agent fusion}: об'єднання інформації між агентами або платформами.
\end{itemize}

Ключова математична тема: усе це є різними способами апроксимації умовного розподілу
\begin{equation}
  p(x_t\mid z_{1:t}^{(1)},z_{1:t}^{(2)},\ldots,z_{1:t}^{(m)}),
\end{equation}
де $x_t$ --- стан, а $z^{(i)}$ --- потік вимірювань $i$-го сенсора.

\section{Калібрування і синхронізація}

Більшість помилок у fusion-системах не є ``помилками нейромережі''. Вони походять з невідомих або змінних параметрів:
\begin{equation}
  T_{ab}\in SE(3),\qquad \Delta t_{ab}\in\R,\qquad K, D, b, \sigma.
\end{equation}
Тут $T_{ab}$ --- extrinsic transform між сенсорами, $\Delta t_{ab}$ --- часовий offset, $K,D$ --- camera intrinsics/distortion, $b$ --- bias, $\sigma$ --- noise scale. Український корпус має підкреслювати: без калібрування злиття сенсорів створює впевнені, але хибні оцінки.

\section{Подієві камери}

Event camera не повертає стандартний кадр $I(x,y,t)$, а події
\begin{equation}
  e_k=(x_k,y_k,t_k,p_k),
\end{equation}
коли лог-інтенсивність у пікселі змінюється на поріг. Це створює асинхронний потік з дуже високою часовою роздільністю і високим динамічним діапазоном. Математично це змушує перейти від ``кадрів'' до подій, потоків, часових поверхонь, накопичувальних вікон і неперервно-часового оцінювання.

\section{Практичний міст до ESKF і factor graphs}

Event/Vision/IMU odometry природно зводиться до вже наявних модулів:
\begin{itemize}[leftmargin=*]
\item кватерніони та $SE(3)$ для орієнтації й пози;
\item ESKF для швидкого рекурсивного оцінювання;
\item factor graphs для batch/smoothing задач;
\item робастні залишки для outliers;
\item NIS/NEES для статистичної перевірки.
\end{itemize}

Отже, цей блок не витісняє Kalman/Lie матеріал, а замикає його на сучасні perception stack задачі.
